\subsubsection*{3.2.7.}
设 $W$ 是 $d$ 维可微流形，设
\[
p:T(W)\longrightarrow W
\]
为其切丛，视为流形。若 $(x_1,\ldots,x_d)$ 为 $W$ 上的局部坐标系，则
\[
(x_1,\ldots,x_d,dx_1,\ldots,dx_d)
\]
为 $T(W)$ 上的局部坐标系。设 $\varepsilon'$ 为此坐标系定义的 $T(W)$ 的定向。它不依赖于 $x_i$ 的选择，因而在整体上定义。

假设 $W$ 可定向。子流形 $W\subset T(W)$（零截面）配上一个定向（分别为其法丛的一个定向），定义 $H_d(T(W))$（分别为 $H_c^d(T(W))$）中的元素 $w$。$W\subset T(W)$ 的自相交为 Euler--Poincaré 示性数：
\[
\int w\wedge w=\chi(W)
\]
这里使用定向 $\varepsilon'$。

\footnote{原印本此句混用了三个维数符号：将维数记为 $x$，将坐标下标写至 $n$，并将上同调类放在 $d$ 次。英文文本在一般陈述中统一使用 $d$；下文的特殊情形使用 $d=n$。}

对消灭圈流形 $V$，其复结构定义定向 $\varepsilon$。它与 $S^n$ 的切丛的定向 $\varepsilon'$ 的关系为
\[
\varepsilon=(-1)^{\frac{n(n-1)}2}\varepsilon'.
\]

\subsubsection*{3.2.8.}
\underline{消灭圈} $\delta$ 在相差符号的意义下良定义。对 $n\ne0$，它是由子流形 $S^n\subset V^\circ$ 定义的 $H_c^n(V^\circ)$ 中的元素。对 $n=0$，空间 $V$ 由两点组成，$\delta$ 是它们之差。使用 $V$ 的复定向，3.2.7 给出
\begin{equation}\tag{3.2.8.1}
(\delta,\delta)
=(-1)^{\frac{n(n-1)}2}\bigl(1+(-1)^n\bigr).
\end{equation}

\subsubsection*{3.2.9.}
令
\[
\widetilde H_c^i(V^\circ)
=\operatorname{Ker}\!\left(
H_c^i(V^\circ)\xrightarrow{\ \operatorname{Tr}\ }\mathbb Z
\right);
\]
当 $i\ne2n$ 时这就是 $H_c^i(V^\circ)$。则
\[
\begin{aligned}
\widetilde H_c^i(V^\circ)&=0 &&\text{当 }i\ne n,\\
\widetilde H_c^n(V^\circ)&=\mathbb Z,
  &&\text{由下列元素生成：}\delta.
\end{aligned}
\]
对偶地，
\[
\begin{aligned}
\widetilde H^i(V)&=0 &&\text{当 }i\ne n,\\
\widetilde H^n(V)&\xrightarrow[\sim]{\ (\,\cdot\,,\delta)\ }\mathbb Z.
\end{aligned}
\]
对 $n\ne0$，$\delta\in H_c^n(V^\circ)$ 在 $H^n(V)$ 中的像在 $n$ 为奇数时为零，在 $n$ 为偶数时可被二整除（且不能再整除）。